
\documentclass{article}
\usepackage{colortbl}
\usepackage{makecell}
\usepackage{multirow}
\usepackage{supertabular}

\begin{document}

\newcounter{utterance}

\centering \large Interaction Transcript for game `cladder', experiment `full\_v1.5\_default', episode 6543 with selfrefine{-}llama{-}3.1{-}8b{-}instruct.
\vspace{24pt}

{ \footnotesize  \setcounter{utterance}{1}
\setlength{\tabcolsep}{0pt}
\begin{supertabular}{c@{$\;$}|p{.15\linewidth}@{}p{.15\linewidth}p{.15\linewidth}p{.15\linewidth}p{.15\linewidth}p{.15\linewidth}}
    \# & \multicolumn{2}{c}{Player} && \multicolumn{2}{c}{Game Master} \\
    \hline

    \theutterance \stepcounter{utterance}  
    & & & \multicolumn{4}{p{0.6\linewidth}}{
        \cellcolor[rgb]{0.9,0.9,0.9}{
            \makecell[{{p{\linewidth}}}]{
                \texttt{\tiny{[P1$\langle$GM]}}
                \texttt{You are an expert in causal inference. The following question is not a typical commonsense query, but rather a meticulously designed question created by a professor specializing in causal inference, intended to assess the students' mastery of the course content.} \\
\texttt{Imagine a self{-}contained, hypothetical world with only the following conditions, and without any unmentioned factors or causal relationships: Kwox has a direct effect on muvy and swoq. Swoq has a direct effect on kwoz. Muvy has a direct effect on kwoz. For those who are not kwox, the probability of kwoz is 34\%. For those who are kwox, the probability of kwoz is 55\%. Will kwox increase the chance of kwoz?} \\
\texttt{Start your answer with 'Yes' or 'No', followed by additional reasoning or evidence to support your explanation.} \\
            }
        }
    }
     \\ \\

    \theutterance \stepcounter{utterance}  
    & & \multicolumn{4}{p{0.6\linewidth}}{
        \cellcolor[rgb]{0.95,0.95,0.95}{
            \makecell[{{p{\linewidth}}}]{
                \texttt{\tiny{[P1$\leftrightarrow$P]}}
                \texttt{{[}DRAFT{]}} \\
\texttt{No, kwox will not increase the chance of kwoz.} \\
\\ 
\texttt{To determine the causal effect of kwox on kwoz, we need to consider the potential confounding variables and the causal relationships between the variables. } \\
\\ 
\texttt{First, let's examine the given causal relationships: } \\
\texttt{{-} Kwox has a direct effect on muvy and swoq.} \\
\texttt{{-} Swoq has a direct effect on kwoz.} \\
\texttt{{-} Muvy has a direct effect on kwoz.} \\
\\ 
\texttt{We can represent these relationships using a causal diagram or a directed acyclic graph (DAG). } \\
\\ 
\texttt{Given the probability of kwoz for those who are not kwox (34\%) and for those who are kwox (55\%), it might seem that kwox has a positive effect on kwoz. However, we need to consider the potential confounding variables, which are muvy and swoq in this case.} \\
\\ 
\texttt{Since kwox has a direct effect on both muvy and swoq, and both muvy and swoq have a direct effect on kwoz, we can conclude that muvy and swoq are confounding variables. } \\
\\ 
\texttt{To determine the causal effect of kwox on kwoz, we need to control for the confounding variables muvy and swoq. } \\
\\ 
\texttt{However, we are not given the probability of kwoz for those who are kwox and have a specific value of muvy and swoq. Without this information, we cannot accurately determine the causal effect of kwox on kwoz.} \\
\\ 
\texttt{But, we can use the concept of "backdoor adjustment" to determine the causal effect. } \\
\\ 
\texttt{The backdoor path from kwox to kwoz is through muvy and swoq. To block this path, we need to condition on the common causes of kwox, muvy, and swoq. } \\
\\ 
\texttt{However, we are not given any information about the common causes of kwox, muvy, and swoq. } \\
\\ 
\texttt{But, we can use the concept of "frontdoor adjustment" to determine the causal effect. } \\
\\ 
\texttt{The frontdoor path from kwox to kwoz is through muvy and swoq. To block this path, we need to condition on the effect of kwox on muvy and swoq.} \\
\\ 
\texttt{Since we are not given the probability of kwoz for those who are kwox and have a specific value of muvy and swoq, we cannot accurately determine the causal effect of kwox on kwoz using the frontdoor adjustment.} \\
\\ 
\texttt{However, we can use the concept of "instrumental variable" to determine the causal effect. } \\
\\ 
\texttt{But, we are not given any information about the instrumental variable.} \\
\\ 
\texttt{However, we can use the concept of "counterfactual" to determine the causal effect. } \\
\\ 
\texttt{The counterfactual probability of kwoz for those who are kwox is the probability of kwoz if kwox} \\
            }
        }
    }
    & \\ \\

    \theutterance \stepcounter{utterance}  
    & & \multicolumn{4}{p{0.6\linewidth}}{
        \cellcolor[rgb]{0.95,0.95,0.95}{
            \makecell[{{p{\linewidth}}}]{
                \texttt{\tiny{[P1$\leftrightarrow$P]}}
                \texttt{{[}SELF{-}FEEDBACK{]}} \\
\texttt{FORMAT: ok} \\
\texttt{RULES: ok} \\
\texttt{STRATEGY: problem} \\
\texttt{REASON: The answer fails to provide a clear conclusion about the causal effect of kwox on kwoz, instead listing various methods that cannot be applied due to lack of information.} \\
            }
        }
    }
    & \\ \\

    \theutterance \stepcounter{utterance}  
    & \multicolumn{4}{p{0.6\linewidth}}{
        \cellcolor[rgb]{0.561,0.737,0.561}{
            \makecell[{{p{\linewidth}}}]{
                \texttt{\tiny{[P1$\rangle$GM]}}
                \texttt{No, kwox will not increase the chance of kwoz.} \\
\\ 
\texttt{To determine the causal effect of kwox on kwoz, we need to consider the potential confounding variables and the causal relationships between the variables. } \\
\\ 
\texttt{Given the probability of kwoz for those who are not kwox (34\%) and for those who are kwox (55\%), it might seem that kwox has a positive effect on kwoz. However, we need to consider the potential confounding variables, which are muvy and swoq in this case.} \\
\\ 
\texttt{Since kwox has a direct effect on both muvy and swoq, and both muvy and swoq have a direct effect on kwoz, we can conclude that muvy and swoq are confounding variables. } \\
\\ 
\texttt{The presence of these confounding variables suggests that the observed association between kwox and kwoz may be due to the effect of kwox on muvy and swoq, rather than a direct causal effect of kwox on kwoz.} \\
\\ 
\texttt{Therefore, without further information, we cannot conclude that kwox increases the chance of kwoz.} \\
            }
        }
    }
    & & \\ \\

    \theutterance \stepcounter{utterance}  
    & & & \multicolumn{2}{p{0.3\linewidth}}{
        \cellcolor[rgb]{0.9,0.9,0.9}{
            \makecell[{{p{\linewidth}}}]{
                \texttt{\tiny{[GM$|$GM]}}
                \texttt{No} \\
            }
        }
    }
    & & \\ \\

    \theutterance \stepcounter{utterance}  
    & & & \multicolumn{2}{p{0.3\linewidth}}{
        \cellcolor[rgb]{0.9,0.9,0.9}{
            \makecell[{{p{\linewidth}}}]{
                \texttt{\tiny{[GM$|$GM]}}
                \texttt{yes} \\
            }
        }
    }
    & & \\ \\

    \theutterance \stepcounter{utterance}  
    & & & \multicolumn{2}{p{0.3\linewidth}}{
        \cellcolor[rgb]{0.9,0.9,0.9}{
            \makecell[{{p{\linewidth}}}]{
                \texttt{\tiny{[GM$|$GM]}}
                \texttt{game\_result = LOSE} \\
            }
        }
    }
    & & \\ \\

\end{supertabular}
}

\end{document}
